
This paper investigates whether the identity of the training source dataset --- independent of training volume and format --- determines the pass@1 performance of a code language model after supervised fine-tuning (SFT). We train DeepSeek-Coder-1.3B-Base on four isolated conditions (HumanEval-only, MBPP-only, LeetCode-only, and Equal-mix) under token-budget equalization, deduplication against test benchmarks, and standardized supervision format, then evaluate on HumanEval+ and MBPP+ using EvalPlus. A one-way ANOVA over HumanEval+ pass@1 at 1.3B scale yields F=11.37, p=0.020, with a maximum pairwise contrast of 31.9 percentage points between HumanEval-only (mean 35.0\%) and LeetCode-only (mean ~3.1\%). Contrary to the symmetric specialization hypothesis, HumanEval-only SFT achieves the highest pass@1 on both HumanEval+ (~35.9\%) and MBPP+ (~51.9\%), while MBPP-only SFT underperforms on MBPP+ relative to HumanEval-only SFT in all three seeds. CodeBERT embedding cosine similarity between training source and test benchmark perfectly predicts the HumanEval+ pass@1 rank across all four conditions (Spearman $\rho$=1.0, permutation test p=0.042, n=10,000; MiniLM $\rho$=0.800, concordant). At 7B scale, the condition rank order is preserved on HumanEval+ with substantially reduced absolute spread (5.5pp versus 31.9pp at 1.3B), though within-condition seed variance collapses to near-deterministic levels, rendering $\eta^2$ an inappropriate cross-scale comparison metric. MBPP+ evaluation is incomplete for the full four-condition matrix; claims about MBPP+ are limited to the two-condition comparison available from H-M1. These findings suggest that embedding-space alignment between training source and evaluation benchmark is a measurable predictor of SFT performance rank, and that algorithmic function-completion training confers broader cross-benchmark generalization than utility-script training at limited fine-tuning scale.

